\subsection{10.32: a logarithmic bound for symmetric-group power words}
\label{cand:10.32}
\begin{proposition}
Let \(r,s\ne0\), with at least one odd, and
\(m=\prod_{\ell\mid rs}\ell\), the product of the distinct prime divisors
of \(|rs|\). The word \(x^ry^s\) is universal on \(S_n\) whenever
\[
 n\ge\max\{1126,16\log m\}.
\]
More generally, every real \(c>8\) admits \(N_0(c)\) for which
\(n\ge\max\{N_0(c),c\log m\}\) suffices.
\end{proposition}
\begin{proof}
We use two permutation-covering results. First, if two classes in
\(S_d\) have \(h,k\) cycles, counting fixed points, and
\[
 h+k+1\le L\le d,\qquad L\equiv h+k-1\pmod2,
\]
some representatives have product an \(L\)-cycle.
This is Boccara's positive-genus three-partition result in the form
reproved by Song--Xu--Ye \cite[Propositions~1.4 and~2.1]{song-xu-ye}.
Second, if \(u,v\ge4\) and
\(\lfloor3n/4\rfloor+1\le u+v\le n\), every even permutation is
a product of two permutations each consisting of a disjoint
\(u\)-cycle and \(v\)-cycle \cite[Lemma~5]{brenner-evans-silberger}.

Suppose an odd \(q\) satisfies \(3n/8+1<q\le n/2\).
Every odd \(g\in S_n\) is a product of a \(2q\)-cycle and two disjoint
\(q\)-cycles. For a transposition \(g\), choose the two \(q\)-cycles
with its letters in different supports; multiplying their inverse
by \(g\) merges them.
Otherwise let \(s_0\) be the support size, \(t\) its number of nontrivial
cycles, and \(e=s_0-t\), an odd integer at least three.
Restrict to an invariant subset of size \(d=\max\{s_0,2q\}\).
The class of \(g\) there has \(h=t+d-s_0\) cycles; the class of two
\(q\)-cycles has \(k=d-2q+2\).
If \(s_0\le2q\), then \(h+k+1=2q-e+3\le2q\).
Otherwise, using \(t\le s_0/2\),
\[
 h+k+1=s_0+t-2q+3\le3n/2-2q+3<2q.
\]
Also \(h+k-1=2d-e-2q+1\) is even.
The first covering result, followed by inversion of its second factor
and simultaneous conjugation, gives the assertion. Removing the
unneeded fixed points before applying it is essential.

Now let \(q\ge5\) be prime in that interval with \(q\nmid rs\), and
suppose \(r\) is odd. For odd targets the factors just obtained have
orders \(2q,q\); for even targets the second covering result with
\(u=v=q\) supplies two factors of order \(q\).
Their orders are coprime to \(r,s\), respectively, so taking inverse
exponents modulo those orders gives \(g=x^ry^s\).
If \(s\) is the odd exponent, reverse the argument's factor order
using \(U^sV^r=V^r(V^{-r}UV^r)^s\).
Negative exponents cause no change.

It remains to find a prime in the interval missing from \(m\).
Let \(\vartheta(x)=\sum_{\ell\le x}\log\ell\).
Rosser--Schoenfeld \cite[Theorems~4 and~9]{rosser-schoenfeld} give
\[
 \vartheta(x)>x\left(1-\frac1{2\log x}\right)\quad(x\ge563),
 \qquad \vartheta(y)<1.01624y\quad(y>0).
\]
For \(n\ge1126\), \(\log(n/2)>6\), and \(1.01624<49/48\).
Consequently
\[
 \vartheta(n/2)-\vartheta(3n/8+1)
 >\frac{29n}{384}-\frac{49}{48}>\frac n{16}.
\]
If all primes in the interval divided \(m\), this would contradict
\(\log m\le n/16\). The explicit bound follows.
The same source gives \(\vartheta(x)=x+o(x)\), so the interval's
prime-log sum is \(n/8+o(n)\), eventually exceeding \(n/c\)
for any \(c>8\). This proves the second statement.
\end{proof}

If both exponents are odd, every degree works: every permutation is
a product of two involutions, by reversing each cycle.
An exponent \(1\) or \(-1\) also makes universality immediate.
The stated constants are sufficient, not optimal.
The candidate contribution is this symmetric-group deduction from
the credited covering and prime estimates, rather than new versions
of those imported theorems.
See \artifact{research/10.32-proof.md} and
\artifact{results/10.32-summary.json}.
